\section{Discussion}\label{sec:discussion}

The message of this paper reduces to two axes. First, on multi-tool tool selection the closest prior activation-steering methods (SEKA and AdaSEKA) carry a structural bias that re-amplifies the first tool at the second decoding step, and the bias leaves a specific footprint in data as a rise in \texttt{repeated\_first\_tool\_rate}. Flipping the sign of the rotation on the same subspace and applying it to the query is enough to remove this footprint and to give consistent gains on Qwen2.5-7B-Instruct across MetaTool Subtask4 and $\tau^2$-bench retail. Second, the rotation-based intervention requires no weight training and no prompt expansion, so it sits below fine-tuning in training cost and below prompt scaffolding in token overhead while remaining orthogonal to, and combinable with, KV-cache optimization.

Whether the gain is attributable to the basis or to the operator is separated by the ablations of \S\ref{sec:exp:ablations}. The collapse of stationary K on feature-shuffled and random bases already rules out the ``any direction works'' alternative. The PCA-of-K row will tune the paper's claim between a strong reading (``the catalog ontology is part of the operator'') and a weak one (``low-rank structure is necessary but the catalog is optimal''); either way the random/shuffled collapse is invariant.

The asymmetry between layer-adaptive and query-only is not a weakness but a source of design guidance. Our operator family beats SEKA in both the short-sequence and long-sequence regimes, so a user selects inside the family by the expected tool-count profile of the task. The layer-boundary sweep in Appendix~\ref{app:layer-sweep} (Phase 2.5, coworker track) quantifies the Pareto frontier of $\tau = 1/4$ against its natural alternatives.

Three limitations remain. Our two operators are history-free, so the emitted-state lack identified by Theorem~\ref{thm:no-memory} is not fully resolved; it is stepped past. What an iterative operator that tracks $P_{\mathrm{emitted}}$ would add on top of this family is an open question. The canonical SEKA reproduction on CounterFact failed in our A6000 environment, so the SEKA head-to-head remains conditional on A100 reproduction (Appendix~\ref{app:seka-repro}). Full $\tau^2$-bench multi-turn success and KV-cache compression are out of scope.

Theorem~\ref{thm:no-memory} applies to \emph{any} stationary operator, so ITI, CAA, PASTA, and Focus Directions share the same structural gap. This is a diagnosis, not a critique; but the theory already predicts that those methods will collide with the same repetition mechanism when extended to multi-tool, and a sign-flipped Q-projection can be the default starting tool for such extensions.
